\documentclass[11pt]{article}
\usepackage[margin=1in]{geometry}
\usepackage{amsmath, amssymb, amsthm}
\usepackage{hyperref}

\title{MAT 3150 --- Homework 3 Hints}
\author{}
\date{}

\begin{document}
\maketitle

These are hints, not finished proofs: each one gives the key idea and
leaves the write-up to you.

% ============================================================
\section*{Problem 1: reverse triangle inequality}

\paragraph{Tools.}
\begin{itemize}
  \item Triangle inequality: $|a+b|\le|a|+|b|$ for all $a,b\in\mathbb{R}$.
  \item For $c\ge 0$: \quad $|a|\le c \iff -c\le a\le c$.
\end{itemize}

\paragraph{Hint.}
Write $x$ as a sum that contains $x-y$:
\[
  x = (x-y) + y .
\]
Apply the triangle inequality to this and rearrange. What upper bound do
you get for $|x|-|y|$?

Now do the same thing with the roles of $x$ and $y$ swapped. What do you
get for $|y|-|x|$? Note that $|y-x|=|x-y|$.

\paragraph{Finishing.} You now have $|x|-|y|\le|x-y|$ and
$-(|x|-|y|)\le|x-y|$. Which of the tools above turns that pair into the
single statement you want?

% ============================================================
\section*{Problem 2: a limit from the definition}

\paragraph{The definition.} $a_n\to L$ means: for every $\varepsilon>0$
there is $N\in\mathbb{N}$ such that $|a_n-L|<\varepsilon$ for all
$n\ge N$.

\paragraph{Step 1: guess $L$.} For large $n$ the ``$+2$'' and ``$+5$'' are
negligible. What is $n/(3n)$?

\paragraph{Step 2: scratch work (don't hand this in as is).}
Put $\left|\dfrac{n+2}{3n+5}-L\right|$ over a common denominator. If your
guess is right, the $n$'s in the numerator cancel and you are left with a
\emph{constant} over something that grows with $n$.

Then make it simpler by \emph{making the denominator smaller}. For
example, $3n+5>3n$, so the fraction gets \emph{bigger} if you replace
$3n+5$ by $3n$. You're allowed to overestimate here: you only need an
upper bound that is still less than $\varepsilon$.

Solve ``simple bound $<\varepsilon$'' for $n$. That tells you how to
choose $N$ (use the Archimedean property to justify that such an $N$
exists).

\paragraph{Step 3: the write-up.} Present it in the forward direction:
\begin{quote}
Let $\varepsilon>0$. By the Archimedean property choose $N\in\mathbb{N}$
with $N>\,$\underline{\hspace{1.5cm}}. Then for all $n\ge N$,
\[
  \left|\frac{n+2}{3n+5}-L\right| = \cdots \le \cdots < \varepsilon .
\]
\end{quote}

% ============================================================
\section*{Problem 3: $\infty$ + bounded $=\infty$}

\paragraph{Definitions.}
\begin{itemize}
  \item $(y_n)$ bounded: there is $M>0$ with $|y_n|\le M$ for all $n$,
        so in particular $y_n\ge -M$.
  \item $x_n\to\infty$: for every $K\in\mathbb{R}$ there is $N$ such that
        $x_n>K$ for all $n\ge N$.
  \item You must show that for every $K$ there is $N$ such that
        $x_n+y_n>K$ for all $n\ge N$.
\end{itemize}

\paragraph{Hint.}
Using $y_n\ge -M$, find a lower bound for $x_n+y_n$ in terms of $x_n$ and
$M$ alone.

Given a target $K$, how large does $x_n$ need to be so that this lower
bound is bigger than $K$? The definition of $x_n\to\infty$ works for
\emph{any} threshold, so apply it to that threshold (not to $K$ itself).

% ============================================================
\section*{Problem 4: averages of an increasing sequence}

Let $S_n=x_1+\cdots+x_n$, so $y_n=S_n/n$. You want $y_{n+1}\ge y_n$ for
every $n$.

\paragraph{Approach A: compute the difference.}
Put $y_{n+1}-y_n=\dfrac{S_{n+1}}{n+1}-\dfrac{S_n}{n}$ over the common
denominator $n(n+1)$, and use $S_{n+1}=S_n+x_{n+1}$. The numerator should
simplify to an expression involving only $x_{n+1}$, $n$, and $S_n$.

Then ask: since $(x_n)$ is increasing, how does each of
$x_1,\dots,x_n$ compare to $x_{n+1}$? What does that tell you about $S_n$
compared to $n\,x_{n+1}$?

\paragraph{Approach B: a weighted average.}
Check that
\[
  y_{n+1}=\frac{n\,y_n + x_{n+1}}{n+1},
\]
so $y_{n+1}$ is a weighted average of $y_n$ and $x_{n+1}$. When is a
weighted average of two numbers at least as large as the smaller of the
two numbers? So what do you need to know about how $y_n$ compares with
$x_{n+1}$? (An average of numbers is never bigger than the biggest one.)

\paragraph{Side question.} Is the hypothesis ``positive'' actually needed
anywhere in your argument? (It's fine to not use it; just notice.)

\end{document}
